\subsection{First nonlinear correction from finite projective capacity}
  \label{subsec:nonlinear-projective-correction}

  \subsubsection*{Context and scope}

    Proposition~B.1 establishes the effective Schr\"{o}dinger equation in the form
    \begin{equation}
      \label{eq:schro-full}
      i\hbar_{\mathrm{eff}}\,\partial_t\psi
      \;=\;
      \hat{H}_{\mathrm{eff}}\,\psi \;+\; \mathcal{N}_{\Pi,\mathrm{BI}}[\psi],
    \end{equation}
    where the correction term $\mathcal{N}_{\Pi,\mathrm{BI}} = \mathcal{O}(A_n/A_n^{\max},\,\ell_\chi^2\Delta_g)$
    collects contributions from Born--Infeld saturation and finite projective capacity.
    Under hypothesis~(H4), this term is negligible and the linear equation is recovered.

    The present section relaxes~(H4) partially: we keep $\ell_\chi^2\Delta_g\ll 1$ (continuum
    regime) but allow the local projective occupancy to be non-negligible.
    The goal is to identify the leading-order structure of $\mathcal{N}_{\Pi,\mathrm{BI}}[\psi]$
    in this regime and show that it takes the form of a cubic nonlinearity of
    Gross--Pitaevskii type, with an effective coupling constant determined by the local
    spectral admissible capacity.

  \subsubsection*{Projective occupancy}

    Let $\Sigma(x,t)$ denote the \emph{local spectral admissible capacity}: the effective
    total admissible amplitude available in the projected-stable sector at point $x$ and
    time~$t$.
    Concretely, it is defined by
    \begin{equation}
      \label{eq:Sigma-def}
      \Sigma(x,t) \;:=\; \sum_n w_n(x,t)\,A_n^{\max}(x,t),
    \end{equation}
    where the sum runs over spectrally admissible modes, $w_n(x,t)$ are local spectral
    weights (normalised to unity), and $A_n^{\max} = c_\chi/\sqrt{\lambda_n}$ is the
    Born--Infeld admissibility envelope of mode~$n$.
    The quantity $\Sigma$ carries the dimension of an amplitude; it encodes the local
    ``room'' left in the projected sector before saturation.

    The \emph{projective occupancy ratio}
    \begin{equation}
      \label{eq:eta-def}
      \eta(x,t) \;:=\; \frac{|\psi(x,t)|^2}{\Sigma(x,t)}
    \end{equation}
    measures the fraction of local admissible capacity currently occupied by the projected
    wavefunction.
    By definition, $\eta \in [0,1]$.
    The regime~(H4) of Theorem~B.1 corresponds to $\eta\ll 1$.

  \subsubsection*{Density-dependent correction to the local phase evolution}

    When $\eta$ is no longer negligible, the finite admissible capacity induces a
    density-dependent correction to the local phase evolution.
    The mechanism is the following.
    Each microconfiguration $\omega\in\Pi^{-1}(x)$ contributes to the projected amplitude
    $\psi(x,t)$ with a phase weight determined by its substrate relaxation rate.
    At low occupancy, these contributions are independent.
    At higher occupancy, the constraint that the total projected amplitude must remain
    within the admissible envelope introduces a correlation: configurations are
    ``competing'' for the available admissible capacity.
    This competition feeds back on the effective local chemical potential.

    More precisely, the effective local phase-evolution rate acquires a correction of the
    form
    \begin{equation}
      \label{eq:mu-expansion}
      \mu_{\mathrm{eff}}(\eta) \;=\; \mu_0 \;+\; \gamma\,\eta \;+\; \mathcal{O}(\eta^2),
    \end{equation}
    where $\mu_0$ is the dilute-regime value (determined by the spectral gap) and
    $\gamma$ is a saturation coefficient with the dimension of an energy.
    The coefficient $\gamma$ encodes the first-order response of the admissible
    capacity constraint to variations of the local occupancy.
    In principle, it can be expressed in terms of the local derivative of the
    Born--Infeld saturation functional with respect to the projected density,
    but its explicit form depends on the detailed microstructure of the preimage
    dynamics and is left as an open derivation.
    Substituting the definition~\eqref{eq:eta-def} of $\eta$ into~\eqref{eq:mu-expansion},
    the correction term contributes to the equation of motion as
    $\mu_{\mathrm{eff}}(\eta)\,\psi = \mu_0\,\psi + (\gamma/\Sigma)\,|\psi|^2\psi +
    \mathcal{O}(|\psi|^4\psi)$.

  \subsubsection*{Effective nonlinear equation}

    Incorporating this correction into~\eqref{eq:schro-full}, the projected dynamics in
    the finite-occupancy regime takes the form

    \begin{equation}
      \label{eq:NLS-cosmo}
      \boxed{
        i\hbar_{\mathrm{eff}}\,\partial_t\psi
        \;=\;
        -\kappa\,\Delta_g\,\psi
        \;+\; V_{\mathrm{eff}}(x,t)\,\psi
        \;+\; g(x,t)\,|\psi|^2\,\psi
        \;+\; \mathcal{O}(|\psi|^4\psi),
      }
    \end{equation}
    where $\kappa = \hbar_{\mathrm{eff}}^2/(2m_{\mathrm{eff}})$ is the dispersion coefficient
    (inherited from §B.13) and the \emph{effective nonlinear coupling} is
    \begin{equation}
      \label{eq:g-eff}
      g(x,t) \;:=\; \frac{\gamma}{\Sigma(x,t)}.
    \end{equation}

    Equation~\eqref{eq:NLS-cosmo} has the form of the Gross--Pitaevskii / nonlinear
    Schr\"{o}dinger equation, but with a fundamental difference in interpretation:
    the coupling $g(x,t)$ is \emph{not postulated} but \emph{derived} from the finite
    local spectral admissible capacity.
    It varies in space and time through $\Sigma(x,t)$, reflecting the local structure of
    the admissible spectral sector.

    \begin{remark}[Structural interpretation of the cubic term]
      \label{rem:cubic-interpretation}
      In this framework, the term $g\,|\psi|^2\psi$ does not represent a literal
      self-interaction of the wavefunction with itself.
      It encodes the \emph{self-rétroaction of projective multiplicity}: the density
      $|\psi|^2$ measures how densely the local projected image is occupied by
      preimage configurations, and this density modifies the effective phase-evolution
      rate via the admissible capacity constraint.
    \end{remark}

    \begin{remark}[Sign of $g$ and focusing/defocusing regimes]
      \label{rem:sign-g}
      The sign of $\gamma$ depends on the dominant correction mechanism.
      If the saturation constraint acts as a repulsive pressure (as in the standard
      Born--Infeld saturating regime), one expects $\gamma > 0$, giving $g > 0$: a
      \emph{defocusing} nonlinearity that further spreads the wavefunction.
      In regimes where the competition between preimages leads to \emph{constructive
      mutual reinforcement} of phase-coherent configurations, one can have $\gamma < 0$
      and $g < 0$: a \emph{focusing} nonlinearity capable of balancing dispersive
      spreading.
      The physically relevant regime for localised structures corresponds to the focusing
      case.
    \end{remark}

  \subsubsection*{Solitonic self-coherent projected states}

    In the focusing regime ($g < 0$) and in one spatial dimension, equation~\eqref{eq:NLS-cosmo}
    admits stationary localised solutions of the form
    \begin{equation}
      \label{eq:soliton-ansatz}
      \psi(x,t) \;=\; \phi(x)\,e^{-i\omega t/\hbar_{\mathrm{eff}}},
    \end{equation}
    where $\phi$ satisfies the nonlinear eigenvalue problem
    \begin{equation}
      \label{eq:soliton-ode}
      -\kappa\,\phi'' \;+\; V_{\mathrm{eff}}\,\phi \;+\; g\,|\phi|^2\,\phi
      \;=\; \omega\,\phi.
    \end{equation}
    In the homogeneous case ($V_{\mathrm{eff}} = 0$, $g$ constant), the exact solution is
    \begin{equation}
      \label{eq:soliton-exact}
      \phi(x) \;\propto\; \operatorname{sech}\!\!\left(\frac{x - x_0}{L}\right),
      \qquad L = \sqrt{\frac{2\kappa}{|\omega|}},
    \end{equation}
    with amplitude and width determined jointly by $\kappa$, $|g|$, and $\omega$.

    Within the Cosmochrony framework, these solutions receive a specific interpretation:

    \begin{proposition}[Solitonic states as self-coherent projective classes]
      \label{prop:soliton-cosmo}
      A solitonic solution of the form~\eqref{eq:soliton-ansatz}--\eqref{eq:soliton-exact}
      corresponds to a class of relational configurations
      $\{\omega \in \Omega : \Pi(\omega) \approx \phi\}$ that is \emph{projectivement
      auto-cohérente}: the tendency of substrate modes to dephase and disperse (the term
      $-\kappa\Delta_g\psi$) is exactly balanced by the saturation-induced localisation
      (the term $g|\psi|^2\psi$).

      Equivalently, these are configurations in which the dispersive divergence of
      preimage phases is compensated by the projective multiplicity feedback, so that the
      projected image remains stationary in shape.
    \end{proposition}

    The dispersive term $-\kappa\Delta_g\psi$ arises, as shown in~§B.13, from the
    divergence of projected mode phases across the preimage $\Pi^{-1}(x)$.
    The nonlinear term $g|\psi|^2\psi$ is the leading correction from the finite
    admissible capacity, which constrains how many independent phase contributions can
    accumulate.
    The soliton is the configuration where these two effects exactly cancel.

  \subsubsection*{Regime diagram}

    The three qualitative regimes are summarised as follows:
    \begin{itemize}
      \item $\eta \ll 1$ (dilute projective regime): $g\,|\psi|^2\psi \approx 0$;
      standard linear Schr\"{o}dinger equation; no localisation beyond external potential.
      \item $\eta \lesssim 1$ (moderate occupancy): $g\,|\psi|^2\psi$ is the leading
      correction; Gross--Pitaevskii-type dynamics; possible soliton formation in the
      focusing case.
      \item $\eta \to 1$ (near-saturated regime): higher-order terms in $\mathcal{O}(\eta^2)$
      become relevant; the effective description departs from exact unitary evolution
      at the level of the projected dynamics, signalling the breakdown of the
      low-occupancy approximation and the onset of strongly constrained projective
      regimes; the
      effective description approaches the boundary of projectability.
    \end{itemize}

  \subsubsection*{Structural status}

    Three points deserve emphasis.

    First, equation~\eqref{eq:NLS-cosmo} is not a modification of quantum mechanics.
    It is the leading-order effective equation in a regime where hypothesis~(H4) partially
    fails.
    In all regimes currently accessible to laboratory experiments, $\eta \ll 1$ and the
    linear Schr\"{o}dinger equation remains valid to all measurable precision.

    Second, the coupling $g(x,t) = \gamma/\Sigma(x,t)$ is not a free parameter.
    It is determined by the local spectral admissible capacity $\Sigma$, which is itself
    fixed by the Born--Infeld admissibility envelope and the local spectral weights.
    The nonlinearity is therefore structurally constrained, not phenomenological.
    In particular, the origin of the coupling $g(x,t)$ can be traced back to the
    Born--Infeld saturation constraint on the substrate amplitudes: the nonlinear
    term is the effective manifestation, at the level of the projected dynamics,
    of the bounded-flux condition $|\partial_t\chi_v| \leq c_\chi$.
    This closes the logical chain from the admissibility bound (Section~\ref{subsec:relational-projection})
    through the linear Schr\"{o}dinger emergence (§B.12) to the present nonlinear correction.

    Third, the equation~\eqref{eq:NLS-cosmo} with the identification~\eqref{eq:g-eff}
    provides a natural bridge to collective quantum phenomena:
    in a regime where many particles share the same projected spatial mode (Bose--Einstein
    condensation), the effective coupling $g(x,t)$ encodes the collective projective
    pressure, making contact with the Gross--Pitaevskii equation of condensate physics
    without postulating inter-particle interactions as fundamental.
    These connections are developed further in the companion papers on superconductivity
    and collective quantum dynamics~\cite{Beau2026g}.
